Several approaches have been proposed for non-linear IV regression. One approach is 
to correct the bias caused by the hidden confounder by estimating the conditional distribution of the treatment given the instrumental. The first literature take this approach are \citet{Carrasco2007, Darolles2011}, in which the conditional distribution is estimated by a ratio of Nadaraya-Watson density estimates. One advantage of this approach is that it allows us to learn arbitrary flexible models. In DeepIV \citep{Hartford2017}, the structural function is modeled by deep neural networks and the density is also estimated via mixture density networks \citep{Bishop2006}. Despite the rich expressiveness of the model, the learning tends to be unstable, especially when the treatment is high-dimensional, due to ``forbidden regression'' issue.

Another approach is to extend the 2SLS algorithm based on non-linear basis functions. The first literature taking this approach is sieve IV \citep{Newey2003, Chen2018}, which uses a finite number of basis functions explicitly specified. KIV (kernel instrumental variable regression) \citep{Singh2019} extends sieve IV to allow for an infinite number of basis functions using RKHS. Although these methods enjoy good theoretical properties, the complexity of models can be limited, since they use prespecified basis functions. Our method, DFIV, utilizes this approach but adaptively learns flexible basis functions from data. 

The third approach is to use the generalized method of moment (GMM) \citep{Hansen1982} to find a structural function so that the residual and the instrumental is independent. Recently, \citet{Bennett2019} proposes DeepGMM, which solves a variational formulation of GMM within a highly flexible function space. Although this approach can handle high-dimensional treatment variable, the learning might not be as stable as 2SLS approach, since it involves solving a smooth zero-sum game.
